% Generated draft: requires mathematical and source review.
\documentclass[11pt]{article}
\usepackage[margin=25mm]{geometry}
\usepackage{fontspec}
\setmainfont{Latin Modern Roman}
\usepackage{amsmath,amssymb,mathtools,unicode-math}
\setmathfont{Latin Modern Math}
\usepackage{enumitem,needspace,xurl,hyperref}
\setlength{\parindent}{0pt}
\setlength{\parskip}{5pt}
\setlength{\emergencystretch}{3em}
\newcommand{\N}{\mathbb{N}}
\newcommand{\Z}{\mathbb{Z}}
\newcommand{\Q}{\mathbb{Q}}
\newcommand{\R}{\mathbb{R}}
\newcommand{\C}{\mathbb{C}}
\newcommand{\F}{\mathbb{F}}
\newcommand{\A}{\mathbb{A}}
\newcommand{\PP}{\mathbb{P}}
\newcommand{\E}{\mathbb{E}}
\newcommand{\eps}{\varepsilon}
\newcommand{\dd}{\,\mathrm{d}}
\begin{document}
\section*{UnsolvedMath ID: 3100080}
\textbf{Source problem code:} catalogue entry
\subsection{Definitions and mathematical statement}
A $t$-ary word is a word over an alphabet $\Sigma$ with $|\Sigma|=t$. The question is explicitly about \emph{substitution instance counts}: it asks whether there is a threshold $k$ such that, for every $K>k$, matching the uniform-random encounter rate for every $(K-1)$-word forces matching it for every $K$-word.

To expose the requested implication without supplying absent definitions, let $\mathcal P_m$ denote the source's intended class of $m$-words, let $I(P,W)$ denote its intended substitution-instance encounter rate of $P$ in $W$, and let $u_t(P)$ denote the corresponding rate for a uniform random $t$-ary word. These are placeholders for definitions missing from the record. The proposed threshold assertion has the form
\[
\exists k\ \forall K>k\ \forall W,\qquad
\left[\forall P\in\mathcal P_{K-1},\ I(P,W)=u_t(P)\right]
\Longrightarrow
\left[\forall P\in\mathcal P_K,\ I(P,W)=u_t(P)\right].
\]
The quantifiers over $k,K$ are intended to concern word-size parameters, but their exact conventions, the class of host words $W$, and whether $k$ may depend on $t$ are not supplied. ``Same rate'' may be an exact limiting equality or an asymptotic condition for finite-word sequences; no normalization or limiting regime is given. The displayed schema therefore preserves the requested implication, rather than claiming a self-contained theorem. In particular, substitution-instance counts cannot silently be replaced by ordinary consecutive-substring frequencies.
\subsection{Short English statement}
Rorabaugh asks whether substitution-instance counts have a quasirandomness threshold: beyond some size, random-like counts for every pattern of one size would force random-like counts at the next size. The source does not define the counts or their limiting regime.
\subsection{Sources}
\begin{itemize}
\item Local UnsolvedMath registry, permanent ID 3100080, exact source code catalogue entry, attributed to Rorabaugh; status open, classification OPEN-TRIAGE.
\item Joshua Cooper, \emph{Combinatorial Problems I Like}, October 2020 snapshot, Words and Codes, source-order bullet 80, \url{https://people.math.sc.edu/cooper/combprob.html}. Supplied HTML extraction accessed July 29, 2026.
\item The supplied research summary relates the question to Cooper's work on quasirandom words and substitution-instance counts, but provides no precise publication or defining reference. It says that a definitive resolution could not be verified. This is a supplied open/unverified assessment, not an independent current-status claim.
\item Provenance flags status and rights NEEDS\_REVIEW: the 2020 page labels the item open, later resolution was not checked item by item, and public access was verified without establishing redistribution permission.
\end{itemize}
\end{document}
